In modern social media platforms, toxic polarization drives user frustration and permanent churn, threatening brand safety and platform advertising revenue. To combat this, digital platforms increasingly deploy content moderation interventions such as bridging algorithms, which suppress confrontational, cross-cutting interactions that fall into users' ``latitude of rejection.'' Using the Dynamic-BABE (Biased Assimilation \& Behavioral Entrenchment) agent-based model, we examine the systemic consequences of such algorithmic interventions in the context of co-evolving dyadic trust. We discover a critical counter-intuitive phenomenon: the Polarization Paradox. While the bridging algorithm successfully curtails behavioral frustration and drops cumulative user churn from 39.18\% to 0.38\% (Cohen's $d = 23.97$, $p < 0.001$), preserving platform revenue (increasing from \$121.64 to \$199.24), it simultaneously exacerbates systemic ideological division. Specifically, we find that polarization increases from 0.9550 to 0.9904 (Cohen's $d = -1.28$, $p < 0.001$). By suppressively shielding agents from disagreeable outgroup opinions, the algorithm prevents the constructive exposure necessary to cultivate tolerance and build dyadic interpersonal trust. Under fragile trust, the absence of direct engagement induces trust segregation---wherein trust concentrates exclusively within homogeneous echo chambers while outgroup relationships wither. Consequently, the network undergoes social fragmentation and structural entrenchment. Our results demonstrate that platforms face a fundamental trade-off: short-term economic stabilization and customer retention are achieved at the cost of long-term democratic cohesion, suggesting that superficial moderation fails to resolve the deep-seated cognitive and relational dynamics of online polarization.
